% ─────────────────────────────────────────────────────────────────────────────
% SHARED coursework preamble — every course inputs THIS.
%
% PREAMBLE ONLY. No \begin{document}, no course-specific notation.
%
% Course-specific macros live at COURSE level, in
%   content/<course>/reference_docs/<course>_macros.tex
% which is inputted AFTER this file. ONE per course, shared by hw/qz/exam --
% macros under a single kind are unreachable from a sibling kind.
%
% That split is the point: the page style, the answer box, and the problem
% header are identical everywhere, so they are defined once; \conv and \ustep
% mean something in DSP and nothing in a networks course, so they are not.
%
% A per-problem file therefore starts:
%   % ─────────────────────────────────────────────────────────────────────────────
% SHARED coursework preamble — every course inputs THIS.
%
% PREAMBLE ONLY. No \begin{document}, no course-specific notation.
%
% Course-specific macros live at COURSE level, in
%   content/<course>/reference_docs/<course>_macros.tex
% which is inputted AFTER this file. ONE per course, shared by hw/qz/exam --
% macros under a single kind are unreachable from a sibling kind.
%
% That split is the point: the page style, the answer box, and the problem
% header are identical everywhere, so they are defined once; \conv and \ustep
% mean something in DSP and nothing in a networks course, so they are not.
%
% A per-problem file therefore starts:
%   % ─────────────────────────────────────────────────────────────────────────────
% SHARED coursework preamble — every course inputs THIS.
%
% PREAMBLE ONLY. No \begin{document}, no course-specific notation.
%
% Course-specific macros live at COURSE level, in
%   content/<course>/reference_docs/<course>_macros.tex
% which is inputted AFTER this file. ONE per course, shared by hw/qz/exam --
% macros under a single kind are unreachable from a sibling kind.
%
% That split is the point: the page style, the answer box, and the problem
% header are identical everywhere, so they are defined once; \conv and \ustep
% mean something in DSP and nothing in a networks course, so they are not.
%
% A per-problem file therefore starts:
%   \input{../../../../docs/latex/coursework_preamble.tex}
%   \input{../../reference_docs/cesc470_macros.tex}
%   \begin{document}
%
% Build-check one file, or a whole assignment:
%   docs/latex/build_tex.sh content/cesc_470/hw/hw01/p01_five_components.tex
%   docs/latex/build_tex.sh content/cesc_470/hw/hw01
%
% Doctrine: docs/directives/coursework-solutions.md
% ─────────────────────────────────────────────────────────────────────────────

\documentclass[11pt]{article}

\usepackage[margin=1in]{geometry}
\usepackage{amsmath, amssymb, mathtools}
\usepackage{graphicx}
\usepackage{float}
\usepackage{booktabs}
\usepackage{longtable}
\usepackage{enumitem}
\usepackage{siunitx}
\usepackage[dvipsnames]{xcolor}
\usepackage{tcolorbox}
\usepackage{fancyhdr}
\usepackage{caption}
\usepackage{tikz}
\usepackage{pgfplots}
\pgfplotsset{compat=1.18}
\usepackage[colorlinks=true, allcolors=NavyBlue]{hyperref}

\captionsetup{font=small, labelfont=bf}
\setlist[enumerate]{itemsep=2pt, topsep=4pt}
\setlist[itemize]{itemsep=2pt, topsep=4pt}

% --- final-answer box -------------------------------------------------------
% Every problem file ends with its result in one of these. The condensed
% solutions document is assembled by lifting exactly these.
\newtcolorbox{answerbox}[1][]{
  colback=Green!6, colframe=Green!55!black,
  boxrule=0.6pt, arc=2pt, left=6pt, right=6pt, top=4pt, bottom=4pt,
  #1
}

% --- citation to course material --------------------------------------------
% \cite{Module 01, PDF p55}  ->  a small italic parenthetical.
% Solutions cite the SLIDE/PAGE a formula came from so a grader (or a future
% reader) can check the claim against what was actually taught. See the
% directive: every non-obvious step carries one.
\newcommand{\srcref}[1]{{\small\itshape(#1)}}

% --- a step that came from OUTSIDE the course materials ---------------------
% Used only where the course is genuinely silent and the problem asks for
% outside research. Marked visually so it is never mistaken for lecture content.
\newcommand{\extref}[1]{{\small\itshape\color{Sepia}[external: #1]}}

% --- callout for the trap in a problem --------------------------------------
\newtcolorbox{trap}[1][]{
  colback=BurntOrange!7, colframe=BurntOrange!70!black,
  boxrule=0.6pt, arc=2pt, left=6pt, right=6pt, top=4pt, bottom=4pt,
  fonttitle=\bfseries, title=Watch out, #1
}

% --- per-problem header -----------------------------------------------------
% \problemheader{8}{15 pts}{Target clock rate}
% The optional 4-argument form carries a learning outcome where the course
% states one (CESC 410 does; CESC 470's HW1 does not).
\newcommand{\problemheader}[3]{%
  \begin{center}\large\bfseries
    Problem #1 \quad\normalsize\mdseries(#2)\\[2pt]
    \large\bfseries #3
  \end{center}
  \vspace{2pt}\hrule\vspace{8pt}
}
\newcommand{\problemheaderlo}[4]{%
  \begin{center}\large\bfseries
    Problem #1 \quad\normalsize\mdseries(#2, #3)\\[2pt]
    \large\bfseries #4
  \end{center}
  \vspace{2pt}\hrule\vspace{8pt}
}

% Course name in the running head. Defined BEFORE \lhead references it, and
% \renewcommand'd by each course's macros file (inputted after this one).
\providecommand{\coursename}{}

\pagestyle{fancy}
\fancyhf{}
\lhead{\small\coursename}
\rhead{\small\thepage}
\renewcommand{\headrulewidth}{0.3pt}

%   % CESC 470 Computer Architecture — course-specific macros ONLY.
%
% Inputted AFTER docs/latex/coursework_preamble.tex, which carries everything
% shared (answerbox, problemheader, srcref, page style). Put nothing general
% here: if another course would want it, it belongs in the shared preamble.

\renewcommand{\coursename}{CESC 470}

% --- performance-equation notation ------------------------------------------
% The course writes these in words on the slides; these render them compactly
% and, more importantly, IDENTICALLY across every problem file.
\newcommand{\CPUtime}{T_{\mathrm{CPU}}}
\newcommand{\IC}{\mathrm{IC}}                    % instruction count
\newcommand{\CPI}{\mathrm{CPI}}
\newcommand{\cycletime}{t_{\mathrm{cycle}}}
\newcommand{\clockrate}{f_{\mathrm{clk}}}
\newcommand{\cycles}{N_{\mathrm{cycles}}}
\newcommand{\speedup}{S}

% Amdahl's law, in the form the slides state it (Module 01, PDF p86, slide 49).
\newcommand{\amdahl}[3]{#1 + \dfrac{#2}{#3}}

%   \begin{document}
%
% Build-check one file, or a whole assignment:
%   docs/latex/build_tex.sh content/cesc_470/hw/hw01/p01_five_components.tex
%   docs/latex/build_tex.sh content/cesc_470/hw/hw01
%
% Doctrine: docs/directives/coursework-solutions.md
% ─────────────────────────────────────────────────────────────────────────────

\documentclass[11pt]{article}

\usepackage[margin=1in]{geometry}
\usepackage{amsmath, amssymb, mathtools}
\usepackage{graphicx}
\usepackage{float}
\usepackage{booktabs}
\usepackage{longtable}
\usepackage{enumitem}
\usepackage{siunitx}
\usepackage[dvipsnames]{xcolor}
\usepackage{tcolorbox}
\usepackage{fancyhdr}
\usepackage{caption}
\usepackage{tikz}
\usepackage{pgfplots}
\pgfplotsset{compat=1.18}
\usepackage[colorlinks=true, allcolors=NavyBlue]{hyperref}

\captionsetup{font=small, labelfont=bf}
\setlist[enumerate]{itemsep=2pt, topsep=4pt}
\setlist[itemize]{itemsep=2pt, topsep=4pt}

% --- final-answer box -------------------------------------------------------
% Every problem file ends with its result in one of these. The condensed
% solutions document is assembled by lifting exactly these.
\newtcolorbox{answerbox}[1][]{
  colback=Green!6, colframe=Green!55!black,
  boxrule=0.6pt, arc=2pt, left=6pt, right=6pt, top=4pt, bottom=4pt,
  #1
}

% --- citation to course material --------------------------------------------
% \cite{Module 01, PDF p55}  ->  a small italic parenthetical.
% Solutions cite the SLIDE/PAGE a formula came from so a grader (or a future
% reader) can check the claim against what was actually taught. See the
% directive: every non-obvious step carries one.
\newcommand{\srcref}[1]{{\small\itshape(#1)}}

% --- a step that came from OUTSIDE the course materials ---------------------
% Used only where the course is genuinely silent and the problem asks for
% outside research. Marked visually so it is never mistaken for lecture content.
\newcommand{\extref}[1]{{\small\itshape\color{Sepia}[external: #1]}}

% --- callout for the trap in a problem --------------------------------------
\newtcolorbox{trap}[1][]{
  colback=BurntOrange!7, colframe=BurntOrange!70!black,
  boxrule=0.6pt, arc=2pt, left=6pt, right=6pt, top=4pt, bottom=4pt,
  fonttitle=\bfseries, title=Watch out, #1
}

% --- per-problem header -----------------------------------------------------
% \problemheader{8}{15 pts}{Target clock rate}
% The optional 4-argument form carries a learning outcome where the course
% states one (CESC 410 does; CESC 470's HW1 does not).
\newcommand{\problemheader}[3]{%
  \begin{center}\large\bfseries
    Problem #1 \quad\normalsize\mdseries(#2)\\[2pt]
    \large\bfseries #3
  \end{center}
  \vspace{2pt}\hrule\vspace{8pt}
}
\newcommand{\problemheaderlo}[4]{%
  \begin{center}\large\bfseries
    Problem #1 \quad\normalsize\mdseries(#2, #3)\\[2pt]
    \large\bfseries #4
  \end{center}
  \vspace{2pt}\hrule\vspace{8pt}
}

% Course name in the running head. Defined BEFORE \lhead references it, and
% \renewcommand'd by each course's macros file (inputted after this one).
\providecommand{\coursename}{}

\pagestyle{fancy}
\fancyhf{}
\lhead{\small\coursename}
\rhead{\small\thepage}
\renewcommand{\headrulewidth}{0.3pt}

%   % CESC 470 Computer Architecture — course-specific macros ONLY.
%
% Inputted AFTER docs/latex/coursework_preamble.tex, which carries everything
% shared (answerbox, problemheader, srcref, page style). Put nothing general
% here: if another course would want it, it belongs in the shared preamble.

\renewcommand{\coursename}{CESC 470}

% --- performance-equation notation ------------------------------------------
% The course writes these in words on the slides; these render them compactly
% and, more importantly, IDENTICALLY across every problem file.
\newcommand{\CPUtime}{T_{\mathrm{CPU}}}
\newcommand{\IC}{\mathrm{IC}}                    % instruction count
\newcommand{\CPI}{\mathrm{CPI}}
\newcommand{\cycletime}{t_{\mathrm{cycle}}}
\newcommand{\clockrate}{f_{\mathrm{clk}}}
\newcommand{\cycles}{N_{\mathrm{cycles}}}
\newcommand{\speedup}{S}

% Amdahl's law, in the form the slides state it (Module 01, PDF p86, slide 49).
\newcommand{\amdahl}[3]{#1 + \dfrac{#2}{#3}}

%   \begin{document}
%
% Build-check one file, or a whole assignment:
%   docs/latex/build_tex.sh content/cesc_470/hw/hw01/p01_five_components.tex
%   docs/latex/build_tex.sh content/cesc_470/hw/hw01
%
% Doctrine: docs/directives/coursework-solutions.md
% ─────────────────────────────────────────────────────────────────────────────

\documentclass[11pt]{article}

\usepackage[margin=1in]{geometry}
\usepackage{amsmath, amssymb, mathtools}
\usepackage{graphicx}
\usepackage{float}
\usepackage{booktabs}
\usepackage{longtable}
\usepackage{enumitem}
\usepackage{siunitx}
\usepackage[dvipsnames]{xcolor}
\usepackage{tcolorbox}
\usepackage{fancyhdr}
\usepackage{caption}
\usepackage{tikz}
\usepackage{pgfplots}
\pgfplotsset{compat=1.18}
\usepackage[colorlinks=true, allcolors=NavyBlue]{hyperref}

\captionsetup{font=small, labelfont=bf}
\setlist[enumerate]{itemsep=2pt, topsep=4pt}
\setlist[itemize]{itemsep=2pt, topsep=4pt}

% --- final-answer box -------------------------------------------------------
% Every problem file ends with its result in one of these. The condensed
% solutions document is assembled by lifting exactly these.
\newtcolorbox{answerbox}[1][]{
  colback=Green!6, colframe=Green!55!black,
  boxrule=0.6pt, arc=2pt, left=6pt, right=6pt, top=4pt, bottom=4pt,
  #1
}

% --- citation to course material --------------------------------------------
% \cite{Module 01, PDF p55}  ->  a small italic parenthetical.
% Solutions cite the SLIDE/PAGE a formula came from so a grader (or a future
% reader) can check the claim against what was actually taught. See the
% directive: every non-obvious step carries one.
\newcommand{\srcref}[1]{{\small\itshape(#1)}}

% --- a step that came from OUTSIDE the course materials ---------------------
% Used only where the course is genuinely silent and the problem asks for
% outside research. Marked visually so it is never mistaken for lecture content.
\newcommand{\extref}[1]{{\small\itshape\color{Sepia}[external: #1]}}

% --- callout for the trap in a problem --------------------------------------
\newtcolorbox{trap}[1][]{
  colback=BurntOrange!7, colframe=BurntOrange!70!black,
  boxrule=0.6pt, arc=2pt, left=6pt, right=6pt, top=4pt, bottom=4pt,
  fonttitle=\bfseries, title=Watch out, #1
}

% --- per-problem header -----------------------------------------------------
% \problemheader{8}{15 pts}{Target clock rate}
% The optional 4-argument form carries a learning outcome where the course
% states one (CESC 410 does; CESC 470's HW1 does not).
\newcommand{\problemheader}[3]{%
  \begin{center}\large\bfseries
    Problem #1 \quad\normalsize\mdseries(#2)\\[2pt]
    \large\bfseries #3
  \end{center}
  \vspace{2pt}\hrule\vspace{8pt}
}
\newcommand{\problemheaderlo}[4]{%
  \begin{center}\large\bfseries
    Problem #1 \quad\normalsize\mdseries(#2, #3)\\[2pt]
    \large\bfseries #4
  \end{center}
  \vspace{2pt}\hrule\vspace{8pt}
}

% Course name in the running head. Defined BEFORE \lhead references it, and
% \renewcommand'd by each course's macros file (inputted after this one).
\providecommand{\coursename}{}

\pagestyle{fancy}
\fancyhf{}
\lhead{\small\coursename}
\rhead{\small\thepage}
\renewcommand{\headrulewidth}{0.3pt}

% CESC 470 Computer Architecture — course-specific macros ONLY.
%
% Inputted AFTER docs/latex/coursework_preamble.tex, which carries everything
% shared (answerbox, problemheader, srcref, page style). Put nothing general
% here: if another course would want it, it belongs in the shared preamble.

\renewcommand{\coursename}{CESC 470}

% --- performance-equation notation ------------------------------------------
% The course writes these in words on the slides; these render them compactly
% and, more importantly, IDENTICALLY across every problem file.
\newcommand{\CPUtime}{T_{\mathrm{CPU}}}
\newcommand{\IC}{\mathrm{IC}}                    % instruction count
\newcommand{\CPI}{\mathrm{CPI}}
\newcommand{\cycletime}{t_{\mathrm{cycle}}}
\newcommand{\clockrate}{f_{\mathrm{clk}}}
\newcommand{\cycles}{N_{\mathrm{cycles}}}
\newcommand{\speedup}{S}

% Amdahl's law, in the form the slides state it (Module 01, PDF p86, slide 49).
\newcommand{\amdahl}[3]{#1 + \dfrac{#2}{#3}}


\begin{document}

\problemheader{6}{3 pts}{Three components included in an ISA}

\subsection*{Problem statement}

List three important components included in an ISA.

\subsection*{Approach}

The course gives an explicit six-item list \srcref{Module 01, PDF p23, slide 11},
introduced as ``ISA Includes''. Three points for three items, so one mark each —
the task is to pick three and say what each contributes. Answering from the
slide rather than from memory matters here: the list is specific, and a
plausible-sounding item that is not on it (``clock rate'', ``cache size'') is an
\emph{organization} property, not an ISA one.

\subsection*{Work}

The complete list from the slide, for reference:

\begin{center}
\begin{tabular}{ll}
  \toprule
  \multicolumn{2}{l}{\textbf{ISA includes} \srcref{PDF p23, slide 11}} \\
  \midrule
  1 & Organization of storage \\
  2 & Data types \\
  3 & Encoding and representing instructions \\
  4 & Instruction set (i.e., opcodes) \\
  5 & Modes of addressing data items/instructions \\
  6 & Program visible exception handling \\
  \bottomrule
\end{tabular}
\end{center}

Three, with what each one fixes:

\paragraph{1. Instruction set (the opcodes).}
The set of operations the machine can perform, and the binary opcode that names
each. This is the most visible part of the ISA — it is what a compiler targets
and what the assembly language spells out, since ``the computer's assembly
language embodies its ISA'' \srcref{PDF p29, slide 15}.

\paragraph{2. Organization of storage — the register set.}
How many programmer-visible registers exist, how wide they are, and what each is
for. The review slide lists ``the computer's register set'' among the ISA's
concerns \srcref{PDF p29}. This is architectural because a program names
registers explicitly; adding one changes what programs can be written.

\paragraph{3. Modes of addressing data items and instructions.}
How an instruction specifies \emph{where} its operand lives — an immediate
constant, a register, a memory address computed from a base and offset, and so
on. Also listed among the ISA's concerns on PDF p29.

Any three of the six are acceptable; the other three are \textbf{data types}
(what widths and formats the machine operates on — integers of various sizes,
floating point), \textbf{encoding and representing instructions} (how the fields
are packed into the instruction word — the constraint behind Problem 3's
$2^{32}$ bound), and \textbf{program-visible exception handling}.

\begin{trap}
Everything in the ISA is \textbf{visible to the program}. Note the qualifier the
slide puts on the last item: ``\emph{program visible} exception handling''.
Caches, pipeline depth, clock rate, and branch predictors are \textbf{not} in
the ISA — they are organization \srcref{PDF p31, slide 16}. Listing one of
those answers a different question.
\end{trap}

\subsection*{Check}

Each of the three passes the Problem 4 visibility test — a program can detect
all of them, because it names opcodes, names registers, and encodes addressing
modes directly in its instructions. Change any one and existing binaries stop
working, which is the defining property of an architectural feature. Consistent
with the ISA definition on PDF p23: ``all the information necessary to write a
machine language program''.

\subsection*{Answer}

\begin{answerbox}
Any three of the six the course lists \srcref{PDF p23}; for example:

\begin{enumerate}[nosep]
  \item \textbf{Instruction set (opcodes)} — the operations the machine can
    perform and their binary encodings.
  \item \textbf{Organization of storage / the register set} — how many
    programmer-visible registers there are, their width and purpose.
  \item \textbf{Addressing modes} — how an instruction specifies where its
    operand is (immediate, register, base $+$ offset, \ldots).
\end{enumerate}

The remaining three: \textbf{data types}, \textbf{instruction encoding}, and
\textbf{program-visible exception handling}.

\medskip
All are \emph{programmer-visible}. Caches, pipeline depth, and clock rate are
\textbf{organization}, not ISA.
\end{answerbox}

\end{document}
